\section*{ID6550: What Maps Contain Vs. What Your Spherical Vortex Model Uses: f(θ)}
\paragraph{Metadata.}
PiID: 2827892125077 | RTEID: RTE:P25435.1046:T4.2101 | UUID: f848755b-f1be-515d-9d4d-a394c8d1f989

\paragraph{Canonical Statement.}
3. **Background flow** - On Earth the background is rigid rotation (constant \textbackslash{}(\textbackslash{}Omega\textbackslash{})) → Coriolis parameter \textbackslash{}(f(\textbackslash{}theta)=2\textbackslash{}Omega\textbackslash{}sin\textbackslash{}phi\textbackslash{}). In cosmology the background is **expansion** \textbackslash{}(H(t)\textbackslash{}) (Hubble flow), not rigid rotation. There’s no universal Coriolis from expansion; instead you get tidal shears and peculiar velocities induced locally by inhomogeneous gravity. (So replace “Coriolis term” in your equations with tidal shear + expansion terms to adapt the model.)

\paragraph{Canonical Equation.}
\[
f(\theta)=2\Omega\sin\phi
\]

\paragraph{Dictionary.}
Relation: f(θ)=2Ωsinφ

\paragraph{Encyclopedia.}
What Maps Contain Vs. What Your Spherical Vortex Model Uses: f(θ) is an algebraic definition, written f(θ)=2Ωsinφ. It is an algebraic definition expressing the quantity directly in terms of the variables on its right-hand side. It is built from Ω, φ, defining f(θ). Within the Trawin Zero Point registry it serves as one element of the framework's mathematical narrative.
